\section{Análisis Orbital y de Telecomunicaciones de la Constelación Starlink}

	\subsection*{Contexto Físico y Objetivo}
	
	La constelación \textbf{Starlink} de SpaceX representa el mayor despliegue de satélites de órbita terrestre baja (LEO) de la historia, diseñados para proporcionar conectividad global de banda ancha y baja latencia. Tras el despliegue inicial desde el vehículo de lanzamiento Falcon 9, los satélites utilizan propulsores de efecto Hall (iónicos) de argón/criptón para elevar su órbita (orbit-raising) hasta su altitud operativa final (típicamente entre $540\text{ km}$ y $550\text{ km}$). Durante esta fase crítica, la telemetría continua y el análisis del entorno atmosférico son fundamentales para garantizar la supervivencia del satélite y la calidad del enlace de comunicaciones (bandas Ku y Ka).
	
	El objetivo de esta práctica es desarrollar un \textbf{dashboard interactivo} que procese y analice los datos de telemetría simulados de un satélite Starlink durante una maniobra de ascenso de $120\text{ minutos}$. Se requiere que implementéis herramientas de análisis numérico para evaluar la posición, velocidad, entorno aerodinámico y capacidad de transferencia de datos del satélite. Todo el análisis debe realizarse bajo el paradigma de \textbf{programación vectorizada}, demostrando un dominio absoluto de la biblioteca computacional moderna.
	
	\subsection*{Datos de Partida}
	
	Copiad y pegad el siguiente bloque de código en vuestro script. Contiene los datos simulados de telemetría (vectores de tamaño $N=1000$) y las matrices de calibración.
	\begin{lstlisting}[language=Python]
		import numpy as np
		
		# Semilla para reproducibilidad
		np.random.seed(42)
		
		# Tiempo de la mision [minutos]
		t = np.linspace(0, 120, 1000)
		
		# Altitud registrada por telemetria [km] (con ligeras oscilaciones)
		h = 250 + 1.2 * t + 5 * np.sin(0.1 * t)
		
		# Potencia de senal recibida en la estacion terrestre (Gateway) [dBm]
		P_rx = -40 - 0.15 * t + 1.5 * np.random.randn(1000)
		
		# Datos atmosfericos tabulados de la U.S. Standard Atmosphere
		h_tab = np.array([200, 300, 400, 500, 600]) # Altitud [km]
		rho_tab = np.array([2.789e-10, 1.95e-11, 2.82e-12, 5.21e-13, 1.13e-13])
		
		# Matriz de acoplamiento de la antena phased-array y retardos medidos
		A_cal = np.array([[1.20, 0.30, 0.10], 
		[0.20, 1.05, 0.20], 
		[0.10, 0.15, 1.10]])
		b_cal = np.array([15.2, 12.5, 14.8])
\end{lstlisting}
	
	\subsection*{Desarrollo Matemático}
	
	Implementad los siguientes 6 apartados secuenciales utilizando \textbf{exclusivamente funciones de NumPy}. El uso de bucles \texttt{for} o \texttt{while} nativos resultará en una fuerte penalización.
	
	\subsection*{1. Resolución de un sistema lineal (Calibración de la Antena)}
	Las estaciones terrestres de Starlink utilizan antenas de arreglo en fase (phased-array). Durante la recepción, existe un acoplamiento electromagnético entre tres sectores principales del arreglo. El retardo real de propagación de la señal $\mathbf{x} = [x_1, x_2, x_3]^T$ se relaciona con el retardo aparente medido $\mathbf{b_{cal}}$ mediante la matriz de interferencia $\mathbf{A_{cal}}$ según la ecuación: $\mathbf{A_{cal}} \mathbf{x} = \mathbf{b_{cal}}$.
	
	\begin{itemize}
		\item \textbf{Se pide:} Resolver el sistema lineal algebraico para hallar el vector de retardos verdaderos $\mathbf{x}$ empleando las rutinas de álgebra lineal de NumPy.
	\end{itemize}
	
	\subsection*{2. Regresión lineal (Modelo de Atenuación de Señal)}
	La potencia de la señal de telemetría $P_{rx}$ sufre una atenuación a lo largo de la maniobra debido al aumento de la distancia y las condiciones ionosféricas.
	
	\begin{itemize}
		\item \textbf{Se pide:} Ajustar los datos dispersos $(t, P_{rx})$ a un modelo lineal $P_{rx}(t) = m \cdot t + c$. Calculad la pendiente $m$ (tasa de degradación en $\text{dBm/min}$) y la ordenada en el origen $c$ mediante un ajuste de mínimos cuadrados. Mostrad la recta de regresión superpuesta a los datos ruidosos.
	\end{itemize}
	
	\subsection*{3. Interpolación 1D (Perfil de Densidad Atmosférica)}
	Para predecir el frenado aerodinámico (drag) en régimen de flujo molecular libre, es necesario conocer la densidad del aire en cualquier instante de la trayectoria.
	
	\begin{itemize}
		\item \textbf{Se pide:} Interpolar el vector de altitud continua $h$ utilizando los datos tabulados discretos \texttt{h\_tab} y \texttt{rho\_tab} para obtener un vector continuo de densidad $\rho(t)$ correspondiente a cada instante de la misión. \textit{Sugerencia: Debido al comportamiento exponencial de la atmósfera, se valorará positivamente interpolar en escala logarítmica.}
	\end{itemize}
	
	\subsection*{4. Resolución de una Ecuación No Lineal (Altitud Crítica)}
	La resistencia atmosférica decae exponencialmente. SpaceX define un umbral de seguridad donde el ``drag'' es considerado insignificante, permitiendo el despliegue total de los paneles solares. Asumiremos que este umbral se alcanza a una altitud exacta de $h_{crit} = 375 \text{ km}$.
	
	\begin{itemize}
		\item \textbf{Se pide:} Encontrar el instante exacto de tiempo $t_{crit}$ en el que el satélite alcanza dicha altitud, resolviendo la raíz de la ecuación no lineal $h(t) - 375 = 0$. Emplead un método vectorizado (ej. localizar el cambio de signo temporal mediante diferencias y aplicar interpolación local lineal).
	\end{itemize}
	
	\subsection*{5. Derivación Numérica 1D (Cinemática del Satélite)}
	A partir de los datos de altitud $h(t)$, debemos caracterizar la tasa de ascenso.
	
	\begin{itemize}
		\item \textbf{Se pide:} Calcular el vector de velocidad radial $v(t) = \frac{dh}{dt}$ (en $\text{m/s}$) y la aceleración radial $a(t) = \frac{d^2h}{dt^2}$ (en $\text{m/s}^2$) utilizando esquemas numéricos de diferencias finitas centrales. Aseguraos de tratar adecuadamente las condiciones en los bordes del array.
	\end{itemize}
	
	\subsection*{6. Integración Numérica 1D (Telemetría de Datos Transmitidos)}
	La capacidad de ancho de banda del satélite depende de la relación señal/ruido. Asumiremos un modelo empírico donde la tasa de bits instantánea en Megabytes por minuto se aproxima exponencialmente a partir de la potencia de señal recibida: $R(t) = 10 \cdot 10^{\frac{P_{rx}(t) + 60}{20}} \quad [\text{MB/min}]$.
	
	\begin{itemize}
		\item \textbf{Se pide:} Calcular el volumen total de datos $V$ (en Gigabytes) transmitidos con éxito a la Tierra durante toda la misión, aplicando integración numérica (ej. Regla del Trapecio compuesto) sobre la curva $R(t)$:
		$$ V = \int_{0}^{t_{final}} R(t) dt $$
	\end{itemize}
	
	\subsection*{Requisitos del Dashboard (GUI Interfaz Gráfica)}
	La práctica debe culminar en una única ventana interactiva programada íntegramente con \texttt{matplotlib} y su módulo \texttt{matplotlib.widgets}.
	
	\begin{enumerate}
		\item \textbf{Panel de Control:} En la zona lateral o inferior de la figura, debéis incluir un widget \texttt{RadioButtons} que permita al usuario cambiar dinámicamente entre las 6 gráficas correspondientes a los 6 cálculos realizados.
		\item \textbf{Interactividad:} Incluir al menos un widget \texttt{Slider}. Por ejemplo, un deslizador que permita modificar el umbral crítico de altitud del apartado 4 (entre $300$ y $450\text{ km}$) y que actualice instantáneamente en la gráfica la posición del punto calculado.
		\item \textbf{Resultados Numéricos:} Utilizad la función \texttt{plt.text} (o equivalente) anclada en un rincón de la figura gráfica para mostrar en texto claro el resultado matemático fundamental del apartado activo (ej. \textit{"Tasa de degradación: -0.15 dBm/min"}, o \textit{"Total de datos: 4.2 GB"}).
	\end{enumerate}
	
	\subsection*{Criterios de Entrega}
	\begin{itemize}
		\item Se entregará un único script \texttt{.py}.
		\item \textbf{Tolerancia Cero a Bucles:} Todo cálculo en arrays debe hacerse mediante \textit{broadcasting} o funciones nativas de \texttt{numpy}. Se restará hasta un 50\% de la nota matemática por usar \texttt{for} o \texttt{while} innecesarios.
		\item El código debe seguir las normas de estilo PEP-8, con un código modularizado (uso de funciones \texttt{def} para separar la matemática de la interfaz gráfica).
		\item Se valorará profundamente la interpretación física de los resultados impresa tanto en los comentarios del código como en la propia interfaz visual.
	\end{itemize}